\chapter{Experiments and Results}
\label{chap:experiments}

% Sources: outputs/*/README.md, experiments/*.py,
%          docs/Design Principles.md

This chapter validates the evaluator against the Integrating paper's~\cite{Benchasattabuse2025Integrating} fidelity results, compares its schedule with schedules found by search under several network conditions, and evaluates cost-quality trade-offs, chain-length scaling, and timing sensitivity. Section~\ref{sec:exp:validation} validates the physical evaluator against the source paper. Section~\ref{sec:exp:comparison} compares the paper schedule with schedules found by search, while Section~\ref{sec:exp:tradeoffs} examines cost-quality trade-offs and alternative objectives. Section~\ref{sec:exp:scaling} studies how the required generation budget changes with chain length, and Section~\ref{sec:exp:timing} evaluates memory dephasing and generation latency. Additional runtime, beam-width, and sweep results are reported in Appendix~\ref{app:results}. Unless the Bridging paper~\cite{Benchasattabuse2026Bridging} is named explicitly, ``the source paper'' refers to the Integrating paper throughout this chapter.

Throughout this chapter and the rest of this thesis, \emph{the optimizer} denotes the software implementing the three search tiers of Chapter~\ref{chap:algorithms} that produces the reported search results.

\section{Experimental Setup}
\label{sec:exp:setup}

\paragraph{Parameters and sweeps}
\label{sec:exp:params}

Unless stated otherwise, all experiments use the beam search with width $w_{\mathrm{beam}}=25$ (Section~\ref{sec:alg:beam}) together with the fixed schedule families of Section~\ref{sec:alg:bruteforce}. They start from the reference configuration $\mathcal{N}_{\mathrm{ref}}$ (Section~\ref{sec:model:network}), notably with $N=10$, $e_d=0.01$, $F_{\min}=0.9$, and $E_{\max}=100$. Each experiment changes only the parameters specified for that experiment; all remaining parameters are held fixed. A complete summary of the experimental parameters and sweep ranges is provided in Table~\ref{tab:params} (Appendix~\ref{app:params}).

\paragraph{Baselines}
\label{sec:exp:baselines}

Three fixed non-adaptive baselines, introduced in Section~\ref{sec:alg:bruteforce}, are used throughout: the \emph{raw chain}, with no purification; the \emph{uniform link-level baseline}, which applies the same copy count and circuit sequence at every hop; and the \emph{paper schedule}, the source paper's hand-picked Fig.\,4 example~\cite{Benchasattabuse2025Integrating}. These are kept as separate baselines because the uniform link-level baseline applies the same purification strategy independently at every hop, whereas the paper schedule combines purification at different spans and stages of the network. The paper itself presents this schedule as a demonstration of scheduling flexibility rather than an optimized design.

\paragraph{Metrics}
\label{sec:exp:metrics}

As defined in Section~\ref{sec:model:objectives}, fidelity $F$ is the Bell-pair fidelity after all purification, backbone, and correction operations, i.e.\ the $w$ component of the error vector at the schedule root; rate is $R=P_{\mathrm{success}}/T_{\mathrm{attempt}}$, where $P_{\mathrm{success}}$ is the probability that a complete schedule attempt succeeds and $T_{\mathrm{attempt}}$ is its deterministic wall-clock duration; and resource cost $C$ counts the number of \Gls{hrgs} copies -- or equivalently \texttt{GenNode} leaves -- consumed per trial. Two comparison modes recur throughout this chapter: \emph{matched-cost}, in which the optimizer is constrained to $C=100$, and \emph{budget-relaxed}, in which it may choose any schedule satisfying $C\le E_{\max}$.

\section{Physics Validation}
\label{sec:exp:validation}

% Source: experiments/fig5_fidelity_vs_noise.py,
%          outputs/headline_experiment_n10/README.md

The evaluator of Section~\ref{sec:alg:evaluator} reproduces the source paper's Fig.\,5~\cite{Benchasattabuse2025Integrating} to four significant figures. At $e_d=0.01$, the raw chain, heralded end-node pumping, and optimistic end-node pumping give fidelities consistent with the corresponding values reported by the paper. This agreement validates the physical model used by the search tiers, but does not establish search optimality. Appendix~\ref{app:validation} gives the corresponding comparison across the swept range of $e_d$.

\section{Optimizer versus the Paper Schedule}
\label{sec:exp:comparison}

This section first compares the search results with the source paper's schedule at the reference configuration $\mathcal{N}_\mathrm{ref}$, then examines the effect of non-ideal inner-qubit noise, outer-photon noise, and heterogeneous per-hop conditions.

\subsection{Headline Comparison}
\label{sec:exp:headline}

% Source: outputs/headline_experiment_n10/README.md,
%          outputs/sweep_network_sensitivity/README.md

At the reference configuration, the optimizer improves on the paper schedule both at the same cost and under the full resource budget. The matched-cost search achieves a higher rate, while the budget-relaxed search reaches a substantially higher rate at lower cost. The result obtained without enabling the same-span pumping operation is kept as a separate comparison. Detailed numerical results are reported in Table~\ref{tab:headline} below.

\begin{table}[ht]
\centering
\caption{Headline comparison at the reference configuration $\mathcal{N}_{\mathrm{ref}}$ ($N=10$, $e_d=0.01$, $F_{\min}=0.9$, $E_{\max}=100$). The paper schedule is the hand-picked Fig.\,4 construction from the Integrating paper~\cite{Benchasattabuse2025Integrating}. Matched-cost search fixes $C=100$; budget-relaxed search allows any schedule with $C\le100$. Rate is $R=P_{\mathrm{success}}/T_{\mathrm{attempt}}$.}
\label{tab:headline}
\small
\begin{tabular}{lcccl}
\hline
Schedule & $C$ & $F$ & Rate & Note \\
\hline
Paper schedule & 100 & 0.93 & 4055.92 & hand-picked \\
Matched-cost & 100 & 0.92 & 4158.14 & +2.5\% rate \\
Budget-relaxed & 60 & 0.91 & 6195.95 & +52.8\% rate, 60\% of cost \\
Budget-relaxed, pumping disabled & 50 & 0.90 & 6713.18 & lower cost still \\
\hline
\end{tabular}
\end{table}

\subsection{Network Generalization}
\label{sec:exp:netgen}

% Source: outputs/sweep_network_sensitivity/README.md

The optimizer was also evaluated under non-ideal inner-qubit errors and different numbers of redundant \Gls{bsm} arms. While the headline comparison uses the source paper's idealized network configuration, the non-ideal configurations introduce physical constraints that can change which schedules remain feasible. At the same generation budget, the paper's schedule falls below the fidelity threshold in both non-ideal configurations tested, whereas rerunning the optimizer under the corresponding physical parameters identifies \emph{feasible schedules}. Figure~\ref{fig:netgen} summarizes the resulting comparisons for the two \Gls{bsm} arm counts; the full numeric comparison is given in Table~\ref{tab:netgen} (Appendix~\ref{app:netgen}).

\begin{figure*}[ht]
    \centering
    \includegraphics[width=0.6\textwidth]{figures/results/fidelity_by_config.png}
    \caption{Fidelity by network configuration and schedule at $N=10$, $e_d=0.01$, $E_{\max}=100$, $\gamma=0$. Hatched bars fall below $F_{\min}=0.9$. The paper schedule is infeasible under both non-ideal configurations; the schedules found by search remain feasible at the same budget.}
    \label{fig:netgen}
\end{figure*}

Rate values for an infeasible schedule are not directly comparable with those of a feasible schedule. Section~\ref{sec:disc:synthesis} discusses the distinction between feasibility and rate more generally.

\subsection{Per-Hop Heterogeneous Networks}
\label{sec:exp:heterogeneous}

% Source: outputs/random_network_adaptation/{30 seeds,100 seeds}/README.md,
%         src/hrgs_scheduler/models/random_network.py

So far, noise has only been varied uniformly across hops (Section~\ref{sec:exp:netgen}). A real deployment, by contrast, will have different link lengths, hardware quality, and error rates from hop to hop. This section evaluates such heterogeneous networks by assigning each hop's length, inner-qubit error rate, and arm count independently at random within the ranges used for this experiment. The search remains non-adaptive; only the supplied network configuration changes.

A five-hop network with one hop purposefully assigned a $12\times$ higher inner-qubit error rate than its neighbours illustrates the benefit of allocating purification resources \emph{unevenly} across a heterogeneous chain, as detailed in Table~\ref{tab:weaklink} (Appendix~\ref{app:weaklink}).

The same tendency holds over a sample of 100 independently generated random $N=6$ networks (seeds 0--99, $e_d=0.01$, $E_{\max}=60$, $F_{\min}=0.9$), reported in full in Appendix~\ref{app:weaklink}. Among them, 80 were found to have at least one feasible uniform link-level candidate. The optimizer improves their rate by a mean of $20.37\%$ and never performs worse in the tested sample (Section~\ref{sec:disc:synthesis}).

\subsection{Noise Sensitivity}
\label{sec:exp:edsweep}

% Source: outputs/sweep_ed_n10/README.md

Sweeping $e_d$ from 0 to 0.02 at $N=10$ shows that the optimizer's rate advantage over the uniform link-level baseline increases with noise. Figure~\ref{fig:edsweep} shows the rate (\ref{fig:edsweep:rate}) and fidelity (\ref{fig:edsweep:fidelity}) curves for the paper baseline, matched-cost optimizer, and budget-relaxed optimizer. At low noise the three variants remain close, with the rate curves separating notably around $e_d=0.005$. The matched-cost curve represents the best cost-$100$ result at each point, not a single fixed schedule structure. Further interpretation is given in Section~\ref{sec:disc:nearoptimal}.

\begin{figure*}[ht]
    \centering
    \begin{subfigure}[b]{0.46\textwidth}
        \centering
        \includegraphics[width=\textwidth]{figures/results/rate_vs_ed.png}
        \caption{Rate vs.\ $e_d$}
        \label{fig:edsweep:rate}
    \end{subfigure}
    \hfill
    \begin{subfigure}[b]{0.46\textwidth}
        \centering
        \includegraphics[width=\textwidth]{figures/results/fidelity_vs_ed.png}
        \caption{Fidelity vs.\ $e_d$}
        \label{fig:edsweep:fidelity}
    \end{subfigure}
    \caption{Rate and fidelity versus depolarizing error $e_d$ at $N=10$ and $E_{\max}=100$ for the paper baseline, the optimizer at matched cost, and the budget-relaxed optimizer. The paper baseline's higher fidelity in the right panel reflects unused surplus above the $F_{\min}=0.9$ floor, purchased with cost and rate that the optimizer instead reallocates.}
    \label{fig:edsweep}
\end{figure*}
% \FloatBarrier

\section{Cost-Quality Trade-offs}
\label{sec:exp:tradeoffs}

\subsection{Pareto Frontiers}
\label{sec:exp:pareto}

% Source: outputs/pareto_frontiers/README.md

With a permissive objective ($F_{\min}=0$, $E_{\max}=150$), the beam search evaluates 180 candidates. Figure~\ref{fig:pareto} shows two projections of the resulting Pareto frontier: fidelity vs.\ cost (\ref{fig:pareto:cost}) and fidelity vs.\ rate (\ref{fig:pareto:rate}). Dominated candidates and the non-dominated frontier under each pair of objectives are distinguished as described in the caption (Section~\ref{sec:alg:prelims}, Pareto frontiers).

In the fidelity-cost projection, the paper schedule lies inside the frontier, which changes little once cost exceeds $C\approx60$. In the fidelity-rate projection, the paper schedule is again inside the frontier, which trades off high fidelity against high rate. The paper schedule and the uniform link-level baseline are marked for reference. Appendix~\ref{app:dags} discusses the schedule \glspl{dag} for the key optima cited here and in Section~\ref{sec:exp:headline}.

\begin{figure*}[ht]
    \centering
    \begin{subfigure}[b]{0.49\textwidth}
        \centering
        \includegraphics[width=\textwidth]{figures/results/fidelity_vs_cost_n10.png}
        \caption{Fidelity vs.\ cost}
        \label{fig:pareto:cost}
    \end{subfigure}
    \hfill
    \begin{subfigure}[b]{0.5\textwidth}
        \centering
        \includegraphics[width=\textwidth]{figures/results/fidelity_vs_rate_n10.png}
        \caption{Fidelity vs.\ rate}
        \label{fig:pareto:rate}
    \end{subfigure}
    \caption{Pareto frontiers at $N=10$, $e_d=0.01$, $E_{\max}=150$, over candidates evaluated under a permissive objective ($F_{\min}=0$). Grey points are dominated; connected points form the non-dominated frontier, with the paper schedule and uniform link-level baseline marked for reference.}
    \label{fig:pareto}
\end{figure*}
% \FloatBarrier

\paragraph{Alternative objectives}
\label{sec:exp:objectives}

Changing the objective from maximizing rate to maximizing fidelity or minimizing cost produces different top-ranked schedules without changing the search algorithm. Maximizing fidelity subject to a rate floor returns a higher-fidelity schedule; minimizing cost with only a fidelity floor can select a substantially lower-rate schedule. Adding an explicit rate floor changes that selection without modifying the search procedure; the complete numerical comparison is given in Appendix~\ref{app:objectives}.

\section{Scaling with Chain Length}
\label{sec:exp:scaling}

\subsection{Budget Scaling}
\label{sec:exp:budget}

% Source: outputs/sweep_min_budget_vs_n/README.md

For $N$ ranging from 10 to 28, the minimum $E_{\max}$ for which a schedule within the searched families reaches $F_{\min}=0.9$ is found by dichotomy, rerunning the beam search at each trial budget. A descriptive power-law fit to the resulting points in log-log space gives
\[
    E_{\max}^{\min} \approx 1.192\,N^{1.666},
\]
with an exponent greater than one over the measured range. The paper's budget is first exceeded at $N=28$, while remaining sufficient at $N=26$. Figure~\ref{fig:budget} shows the measured points together with the paper's linear budget and the fitted power law. These fitted values describe the minimum budget found within the searched families at the chosen beam width, not a global lower bound or an asymptotic result (Section~\ref{sec:model:scope}). Chapter~\ref{chap:discussion} discusses this scope limitation together with the excluded-construction result of Section~\ref{sec:exp:excluded}. Appendix~\ref{app:sweeps} additionally reports the minimum feasible budget versus $e_d$ at fixed $N=10$.

\begin{figure*}[ht]
    \centering
    \includegraphics[width=0.85\textwidth]{figures/results/min_budget_vs_n.png}
    \caption{Minimum feasible $E_{\max}$ to reach $F_{\min}=0.9$ vs.\ $N$ at $e_d=0.01$, compared with the paper's linear $E_{\max}=10N$ budget and the fitted power law $E_{\max}^{\min}\approx1.192\,N^{1.666}$. Solid markers and the solid fit cover the measured range ($N=10$--$28$); the dotted extension is the same fit evaluated beyond that range, not additional data.}
    \label{fig:budget}
\end{figure*}
\FloatBarrier

The related optimality-scope check at $N=18$ is reported in Section~\ref{sec:exp:excluded} and Appendix~\ref{app:results}. There, a hand-constructed schedule using a same-span construction excluded from the search reaches the paper's own budget.

\section{Timing Sensitivity}
\label{sec:exp:timing}

% Source: outputs/sweep_gamma_and_tau_emit/README.md

Memory dephasing and generation latency affect schedules through their timing structure, including the placement of \texttt{HeraldNode}s and \texttt{IdleNode}s (Section~\ref{sec:model:dag}). Sweeping the memory dephasing rate $\gamma$ leaves the raw chain and paper schedule unchanged over the tested range, while heralded end-node pumping degrades as $\gamma$ increases because its sacrificial copies remain in memory during the round-trip heralding delays; the schedules selected by the optimizer remain flat over the tested range because the search favours optimistic schedules. Sweeping the generation latency $\tau_{\mathrm{emit}}$ reduces the rate of every schedule. Detailed timing curves (Figure~\ref{fig:timing}) and the corresponding parameter sweeps are given in Appendix~\ref{app:timing}.